\chapter{Finite-Size Mesoscopic Theory Bridge}
\label{app:finite-size-mesoscopic-theory}

This appendix collects the derivation path that turns LIF/GIF spike generators
into implementable finite-size population equations. It is written as a
technical bridge, not as a replacement for the chapter sequence: Chapter~3
introduces the single-neuron and population objects, Chapter~8 instantiates
them in a paired E/I clustered network, and Chapter~9 uses them for diagnostic
protocols.

\section{Population Labels and Microscopic Spikes}

Let \(P=(a,\alpha)\) denote one population: cell type
\(\alpha\in\{E,I\}\) in cluster \(a\). Its local size is
\(N_P=n_\alpha\). The microscopic spike train of neuron \(i\in P\) is
\begin{equation}
  s_i^P(t)=\sum_k \delta(t-t_{i,k}^P),
  \label{eq:appA-spike-train}
\end{equation}
and the empirical population activity is
\begin{equation}
  A_N^P(t)
  =
  \frac{1}{N_P}\sum_{i=1}^{N_P}s_i^P(t).
  \label{eq:appA-empirical-activity}
\end{equation}
The infinite-population activity will be written \(A^P(t)\). The finite-size
coordinate for this population is
\begin{equation}
  \xi_P=N_P^{-1/2},
  \qquad
  \xi_E=n_E^{-1/2},
  \qquad
  \xi_I=n_I^{-1/2}.
  \label{eq:appA-xi}
\end{equation}
In paired-cluster discussions a bare \(\xi\) usually means \(\xi_E\).

\section{Hazard Description}

A stochastic LIF or GIF neuron becomes a point-process neuron when its state is
mapped to a conditional intensity, with \(dN_i^P(t)\) denoting the infinitesimal
spike count:
\begin{equation}
  \Pr\{dN_i^P(t)=1\mid\mathcal H_t\}
  =
  \lambda_i^P(t\mid\mathcal H_t)\,dt+o(dt).
  \label{eq:appA-hazard-definition}
\end{equation}
For a soft-threshold GIF model one may write schematically
\begin{equation}
  \lambda_i^P(t)
  =
  h_\alpha\!\left(
    U_i^P(t)-\vartheta_i^P(t)
  \right),
  \qquad
  \vartheta_i^P(t)
  =
  \theta_0^\alpha+\sum_{t_{i,k}^P<t}
  \eta_\theta^\alpha(t-t_{i,k}^P),
  \label{eq:appA-gif-hazard}
\end{equation}
with an absolute refractory period represented either by
\(\lambda_i^P(t)=0\) for \(t-\hat t_i^P<\tau_{\mathrm{ref}}\) or by a large
short-lag threshold kernel. The variable \(U_i^P\) is the effective
subthreshold drive, including filtered synaptic current and any voltage
rescaling used by the neuron model.

\section{Age Density}

Let \(\ell=t-\hat t_i^P\) be time since the most recent spike. The
population-level age density \(q^P(\ell,t)\) is normalized by
\begin{equation}
  \int_0^\infty q^P(\ell,t)\,d\ell=1.
  \label{eq:appA-age-normalization}
\end{equation}
If all neurons of the same age share the same effective hazard
\(\lambda^P(\ell,t)\), conservation of neurons gives
\begin{align}
  \partial_t q^P(\ell,t)+\partial_\ell q^P(\ell,t)
  &=
  -\lambda^P(\ell,t)q^P(\ell,t),
  \qquad \ell>0,
  \label{eq:appA-age-pde}
  \\
  q^P(0,t)
  &=
  A^P(t),
  \label{eq:appA-age-boundary}
  \\
  A^P(t)
  &=
  \int_0^\infty
  \lambda^P(\ell,t)q^P(\ell,t)\,d\ell.
  \label{eq:appA-age-activity}
\end{align}
The mass check is immediate:
\[
  \frac{d}{dt}\int_0^\infty q^P(\ell,t)d\ell
  =
  q^P(0,t)-\int_0^\infty\lambda^P(\ell,t)q^P(\ell,t)d\ell
  =
  A^P(t)-A^P(t)=0.
\]

\section{Renewal and Quasi-Renewal Forms}

Given a last spike at time \(t-\ell\), define
\begin{align}
  S^P(\ell\mid t-\ell)
  &=
  \exp\!\left[
    -\int_0^\ell \lambda^P(u,t-\ell+u)\,du
  \right],
  \label{eq:appA-survival}
  \\
  f^P(\ell\mid t-\ell)
  &=
  \lambda^P(\ell,t)S^P(\ell\mid t-\ell).
  \label{eq:appA-isi-density}
\end{align}
The renewal activity equation is
\begin{equation}
  A^P(t)
  =
  \int_0^\infty
  f^P(\ell\mid t-\ell)A^P(t-\ell)\,d\ell
  + A_{\mathrm{init}}^P(t).
  \label{eq:appA-renewal}
\end{equation}
For stationary input the initialization term disappears and
\begin{equation}
  \bar A^P
  =
  \left[\int_0^\infty S^P(\ell)\,d\ell\right]^{-1}.
  \label{eq:appA-stationary-rate}
\end{equation}

For GIF neurons, the exact hazard depends on the full spike history through
threshold and adaptation kernels. A quasi-renewal closure keeps the last spike
exact through \(\ell\) and replaces older spikes by their expected population
history:
\begin{equation}
  \vartheta_{\mathrm{QR}}^P(\ell,t)
  =
  \theta_0^\alpha
  + \eta_\theta^\alpha(\ell)
  + \int_\ell^\infty
  \eta_\theta^\alpha(s)A^P(t-s)\,ds,
  \qquad
  \lambda^P(\ell,t)
  \approx
  h_\alpha\!\left(U^P(t)-\vartheta_{\mathrm{QR}}^P(\ell,t)\right).
  \label{eq:appA-quasi-renewal}
\end{equation}
This is the practical reduction: the neuron model contributes \(h_\alpha\) and
history kernels, while the population equation advances \(q^P\) and \(A^P\).

\section{Finite-Size Expansion}

In a bin of width \(\Delta t\), the empirical activity estimator is
\begin{equation}
  A_{N,\Delta}^P(t)
  =
  \frac{1}{N_P\Delta t}
  \sum_{i=1}^{N_P}
  \int_t^{t+\Delta t}s_i^P(u)\,du.
  \label{eq:appA-binned-activity}
\end{equation}
Conditioned on the population state and ignoring shared input within the bin,
\begin{equation}
  \mathbb E[A_{N,\Delta}^P(t)]=A^P(t),
  \qquad
  \operatorname{Var}[A_{N,\Delta}^P(t)]
  \approx
  \frac{A^P(t)}{N_P\Delta t}.
  \label{eq:appA-binned-moments}
\end{equation}
The corresponding diffusion notation is
\begin{equation}
  A_N^P(t)
  =
  A^P(t)
  +
  \xi_P\sqrt{A^P(t)}\,\zeta^P(t),
  \label{eq:appA-finite-size-activity}
\end{equation}
where \(\zeta^P\) is a shorthand for the activity noise. In renewal theory this
noise need not be white after the refractory structure is resolved; for
clustered-network scaling arguments, the essential point is the prefactor
\(\xi_P=N_P^{-1/2}\).

Filtered rates are
\begin{equation}
  R_N^P(t)=(\epsilon*A_N^P)(t).
  \label{eq:appA-filtered-rate}
\end{equation}
For the exponential kernel, an implementation with bin width \(\Delta\) can use
\begin{equation}
  R_{m+1}^P
  =
  e^{-\Delta/\tau_s}R_m^P
  +
  \left(1-e^{-\Delta/\tau_s}\right)A_{N,\Delta,m}^P.
  \label{eq:appA-filter-recursion}
\end{equation}
For independent stationary Poisson spike trains this gives the covariance
check
\begin{equation}
  C_{R_N^P}(\tau)
  =
  \frac{\bar A^P}{2\tau_s N_P}e^{-|\tau|/\tau_s}.
  \label{eq:appA-filtered-covariance}
\end{equation}
The \(1/N_P\) amplitude is the part used for finite-size scaling; deviations
from the Poisson shape diagnose refractory structure, common input, or
collective dynamics.

\section{Coupled Cluster Populations}

For a paired E/I clustered network, one population equation is written for each
\(P=(a,\alpha)\):
\begin{align}
  \partial_t q_a^\alpha(\ell,t)+\partial_\ell q_a^\alpha(\ell,t)
  &=
  -\lambda_a^\alpha(\ell,t)q_a^\alpha(\ell,t),
  \label{eq:appA-cluster-age-pde}
  \\
  q_a^\alpha(0,t)
  &=
  A_a^\alpha(t),
  \qquad
  A_a^\alpha(t)
  =
  \int_0^\infty
  \lambda_a^\alpha(\ell,t)q_a^\alpha(\ell,t)d\ell.
  \label{eq:appA-cluster-activity}
\end{align}
The drive entering the hazard is obtained from filtered synaptic input. At the
population bookkeeping level,
\begin{equation}
  y_a^\alpha(t)
  =
  \sum_{\beta}
  n_\beta\widehat J^{\alpha\beta}A_a^\beta(t)
  +
  \sum_{b\ne a}\sum_{\beta}
  n_\beta\check J_{ab}^{\alpha\beta}A_b^\beta(t),
  \qquad
  I_a^\alpha=\epsilon*y_a^\alpha.
  \label{eq:appA-cluster-input}
\end{equation}
Chapter~8 supplies the signs, fixed in-degree construction, balance scaling,
and potentiation/depression factors. The finite-size mesoscopic version
replaces \(A_a^\alpha\) by
\begin{equation}
  A_{a,N}^\alpha(t)
  =
  A_a^\alpha(t)
  +
  \xi_\alpha\sqrt{A_a^\alpha(t)}\,\zeta_a^\alpha(t),
  \label{eq:appA-cluster-noisy-activity}
\end{equation}
while frozen inter-cluster disorder is represented separately, for example by
\begin{equation}
  \check J_{ab}^{\alpha\beta}
  =
  \overline J^{\alpha\beta}
  +
  \kappa G_{ab}^{\alpha\beta},
  \qquad a\ne b.
  \label{eq:appA-kappa-decomposition}
\end{equation}
Taking \(n_E,n_I\to\infty\) sends \(\xi_E,\xi_I\to0\). It does not remove
\(\kappa\) unless the network ensemble is scaled to remove inter-cluster
heterogeneity.

\section{Derivation Checks for Implementations}

\begin{enumerate}
  \item Conservation: numerical advection of \(q_a^\alpha(\ell,t)\) should keep
  \(\int q_a^\alpha d\ell=1\) up to discretization error.

  \item Renewal consistency: under stationary input, the simulated population
  rate should match the inverse mean ISI computed from the survival function.

  \item Filtering consistency: filtering binned spikes with
  \eqref{eq:appA-filter-recursion} should match direct convolution with
  \(\epsilon\) when the same bin convention is used.

  \item Finite-size scaling: at fixed input and fixed quenched network draw,
  independent sampling variance should scale as \(1/n_\alpha\), while the
  standard deviation scales as \(1/\sqrt{n_\alpha}\).

  \item Disorder separation: changing random seeds for spike generation probes
  \(\xi\); changing the sampled inter-cluster graph or weights probes
  \(\kappa\). These sweeps should not be merged when interpreting switching.
\end{enumerate}
